% Modello "Tesi" — LaTeX Editor Offline
% Contenuto segnaposto originale: libero da usare e modificare.
\documentclass[12pt,a4paper]{report}

\usepackage[T1]{fontenc}
\usepackage[utf8]{inputenc}
\usepackage[italian]{babel}
\usepackage{amsmath,amssymb}
\usepackage{graphicx}
\usepackage{booktabs}
\usepackage{setspace}
\usepackage[hidelinks]{hyperref}

\onehalfspacing

\newcommand{\tesiTitolo}{Titolo della tesi}
\newcommand{\tesiCandidato}{Nome Cognome}
\newcommand{\tesiCorso}{Corso di laurea di esempio}
\newcommand{\tesiRelatore}{Prof. Nome Relatore}
\newcommand{\tesiAnno}{2025/2026}

\begin{document}

% ---------------------------------------------------------- frontespizio ---
\begin{titlepage}
  \centering
  {\Large Universit{\`a} di esempio\par}
  \vspace{0.5cm}
  {\large \tesiCorso\par}
  \vspace{3cm}
  {\Huge \tesiTitolo\par}
  \vspace{3cm}
  {\large Candidato:\par}
  {\large \tesiCandidato\par}
  \vspace{1cm}
  {\large Relatore:\par}
  {\large \tesiRelatore\par}
  \vfill
  {\large Anno accademico \tesiAnno\par}
\end{titlepage}

% -------------------------------------------------------------- ringraziamenti
\chapter*{Ringraziamenti}
Ringrazio chi ha contribuito alla realizzazione di questo lavoro.
\textit{(Segnaposto: rimuovi il capitolo se non serve.)}

% ------------------------------------------------------------------- indice
\tableofcontents

% ------------------------------------------------------------------- capitoli
\chapter{Introduzione}
Descrivi il problema, la motivazione e la struttura della tesi. Una
citazione di esempio: \cite{esempio2024}.

\chapter{Stato dell'arte}
\section{Lavori correlati}
Rassegna della letteratura, organizzata per temi.

\section{Limiti delle soluzioni esistenti}
Evidenzia il margine di contributo della tesi.

\chapter{Metodologia}
Descrivi il metodo proposto. Un'equazione di esempio:

\begin{equation}
  \label{eq:metodo}
  f(x) = ax^2 + bx + c
\end{equation}

e il riferimento ad essa: \eqref{eq:metodo}.

\chapter{Sperimentazione}
\section{Impostazione sperimentale}
Protocollo, dataset e metriche.

\section{Risultati}
Riporta i risultati con tabelle e figure, per esempio:

\begin{table}[htbp]
  \centering
  \caption{Risultati di esempio.}
  \begin{tabular}{lrr}
    \toprule
    Configurazione & Metrica A & Metrica B \\
    \midrule
    Base           & 0,72      & 0,81      \\
    Proposta       & 0,85      & 0,90      \\
    \bottomrule
  \end{tabular}
\end{table}

\chapter{Conclusioni}
Sintesi del contributo, limiti e sviluppi futuri.

% ------------------------------------------------------------------ bibliografia
\bibliographystyle{plain}
\bibliography{bibliography}

\end{document}
